\section{Introduction}
\label{sec:conductor-introduction}
The multiplication map of a linear series can fail to be surjective
with a scheme structure which changes even when its support does
not. We study this structure for series containing the sections
vanishing on a finite divisor. The method has two parts: a global
conductor theorem identifies the actual cokernel with a finite-algebra
cokernel, and a quotient-algebra description determines the latter
scheme without passing to its radical.

Two results take this comparison beyond reduced contacts and rank-two
quotients. First, the stable generated algebra carries a representation
on the multiplication defect. Its exact-rank strata identify the
conductor quotient as a flag $C\subset Q$ with $C$ acting faithfully
on $Q/C$ (Theorem~\ref{thm:higher-defect-strata}). In codimension two
there are two mechanisms: a rank-three quotient with its scalar line,
or a rank-four quotient which is quadratic over a quadratic algebra
(Theorem~\ref{thm:codimension-two}). The transition between them can
itself be nonreduced.

Second, let $Z$ be the order-$h$ fat point of $\Pj^e$, with $e\ge2$
and $h\ge3$. For every $n\ge2h-1$, take all series containing
$H^0(\mathcal I_Z(n))$ and restricting to a generating $(e+1)$-plane
on $Z$. In symmetric degrees $m\ge h-1$ their failure ideal is
\begin{equation}\label{eq:intro-fat-main}
 \mathcal I_\Delta^{\binom{e+h-1}{e+1}},
\end{equation}
where $\Delta$ is the integral determinant divisor of the map from
the augmentation kernel of the restricted series to
$\mathfrak m/\mathfrak m^2$. This is the full primary ideal, and
all its powers remain primary (Corollary~\ref{cor:global-fat-primary}).
The proof separates a filtered substitution determinant from a
triangular global conductor argument. It gives three-plane results
on every $\C[x,y]/(x,y)^h$ and works in arbitrary embedding dimension;
no diagonal-ideal theorem is used. The global range depends on the
fat-point order, rather than its binomially growing length.


Let $D\subset\Pj^1$ have length $d$, put
$E_{n,D}=H^0(\mathcal O_D(n))$ and $W_D=H^0(\mathcal O(n)(-D))$,
and let $A\subset E_{n,D}$ be a generating $k$-plane. Write
$U_A=\operatorname{res}_D^{-1}(A)$. For $1\le k<d$ our conductor
range is $n\ge2d-k$. In every symmetric degree $m\ge2$ it gives
\begin{equation}\label{eq:intro-conductor-cokernel}
 \operatorname{coker}(\Sym^m\mathcal U_A\to V_{mn})
 \simeq\operatorname{coker}(\Sym^m\mathcal A\to E_{mn,D}).
\end{equation}
The isomorphism holds on the entire relative generating Grassmannian
and after arbitrary base change. The range is sharp uniformly in
$D,A,m$: at $n=2d-k-1$ an explicit monomial series already violates
it in degree two. The exact fibrewise obstruction is the first
cohomology of an evaluation kernel. On a curve of genus $g$, intrinsic
cohomological hypotheses give the same transport; the uniform
sufficient bound is $\deg L\ge2d+2g-1$.

The defect-one structural theorem concerns generating hyperplanes in
an arbitrary finite locally free algebra, not only pencils or powers
of one element. Its failure scheme in all degrees $m\ge2$ is the
space of rank-two quotient algebras with a generating line. The
identification is an isomorphism of functors on arbitrary, including
nonreduced, test schemes. It follows from the conductor of a unital
hyperplane subalgebra and a cyclic presentation of the multiplication
cokernel (Theorem~\ref{thm:universal-hyperplane}).

For $D=\sum_i d_i p_i$ this gives a complete primary classification
of all generating $(d-1)$-planes. The components supported on $2p_i$
have transverse algebra
\[
 G_{d_i}=\C[u,v]/(F_{d_i},F_{d_i+1}),\qquad
 F_0=0,\quad F_1=1,\quad F_{j+1}=uF_j+vF_{j-1};
\]
those supported on $p_i+p_j$ have transverse algebra
$\C[x,y]/(x^{d_i},y^{d_j})$. Their lengths are $\binom{d_i}2$
and $d_id_j$, but their exact nilpotency indices are $2d_i-3$ and
$d_i+d_j-1$. The reduced components are disjoint smooth curves on
the generating open. The moving scheme is smooth and is a generating
open of a projective-line bundle over
$\operatorname{Sym}^2 C\times\operatorname{Sym}^{d-2} C$.
These results apply to global multiplication on $\Pj^1$ for $n\ge d+1$.
They classify the whole contact-hyperplane family for every
multiplicity partition, not a collection of transverse examples.

There is also a complete stable description for generating
three-planes on a reduced divisor. On unit frames its ideal is
\[
 \bigcap_{i<j}(x_i-x_j,y_i-y_j).
\]
The equality, and the analogous equality for every power of the
ideal, use Haiman's diagonal-ideal theorem explicitly. The step
specific to multiplication is an exact total-degree truncation and
the conductor comparison; the radical-ideal theorem is a classical
input, not a new assertion here.

\subsection*{The pencil family and the distinction between fibres and total spaces}
The earlier contact-pencil classification remains part of the article,
now in the improved range. We record its complete statement to fix
the bases of its primary assertions.

\begin{theorem}[Contact-pencil failure schemes]\label{thm:contact-main}
Let $d\ge4$ and $n\ge2d-2$. On the generating open
$G_D^\circ\subset\Gr(2,E_{n,D})$, let $\mathcal J_m$ be the
zeroth Fitting ideal of the global multiplication cokernel.
For $2\le m<d-1$, this ideal is zero: the failure scheme is the
whole parameter space. For $m\ge d-1$ it is a Cartier divisor,
and, on a fixed divisor $D=\sum_i d_i p_i$, its complete primary
decomposition is
\begin{equation}\label{eq:intro-contact-primary}
 \mathcal J_m=
 \bigcap_{i:d_i\ge2}\mathcal I_{T_i}^{\binom{d_i}{2}}
 \ \cap\!\bigcap_{i<j}\mathcal I_{E_{ij}}^{d_i d_j}.
\end{equation}
Here $T_i$ and $E_{ij}$ are failure to separate $2p_i$ and
$p_i+p_j$. These are distinct smooth irreducible divisors; there
are no embedded associated points. The exact local equations,
all coranks and infinitesimal layers are given in
Theorems~\ref{thm:contact-factorization} and~\ref{thm:contact-layers}.
Over a moving multiplicity stratum, primary descent is governed by
Lemma~\ref{lem:etale-primary-descent}.

With $D$ unrestricted, the pencil failure divisor is integral, and
its normalization is the smooth finite incidence marking a scalar
length-two subdivisor. The normalization and its singular-support
test are given in Theorem~\ref{thm:moving-normalization} and
Proposition~\ref{prop:moving-singular-test}.
\end{theorem}

On a unit pencil frame, with $D=(g)$ and ratio $v$, the determinant
is the classical index form
$\Delta(g,v)=\det[1,v,\ldots,v^{d-1}]$. Its relation to the
characteristic polynomial $\chi_v$ is
\[
 \operatorname{disc}(\chi_v)=\operatorname{disc}(g)\Delta(g,v)^2.
\]
The equation of failure is $\Delta=0$, not the discriminant on the
left: the latter also vanishes when $g$ has repeated roots, even if
$v$ generates the algebra. The quotient and index-form constructions
are related to the schemes of generators in \cite{ABHS}. The
Grassmannian cohomology presentation of $G_r$ is likewise classical.
What must be proved here is the exact identification with the
multiplication schemes, including primary fibres and base change.

The dependence of the arguments is
\[
 \begin{gathered}
 \text{evaluation kernel}\ \Longrightarrow\ \text{conductor cokernel},\\
 \text{unital hyperplane conductor}\ \Longrightarrow\ \text{rank-two quotients},\\
 \text{rank-two quotients}\ \Longrightarrow\ \text{primary fibres and smooth total scheme},\\
 \text{index form}\ \Longrightarrow\ \text{pencil primary fibres and normalization},\\
 \text{polarized Cayley--Hamilton + Haiman}\ \Longrightarrow\ \text{three-plane schemes}.
 \end{gathered}
\]
None of these arguments depends on the statistical interpretation,
the quadratic excess classification, or the unrestricted hyperplane
embedded structure. Those developments, including all their proofs
and hypotheses, are retained in the full complementary manuscript.
The present family is not identified with the unrestricted ambient
Grassmannian. Section~\ref{sec:finite-prior-comparison} gives the
precise classical inputs and the remaining limitation in the
nearest-source comparison.
